% Options for packages loaded elsewhere
% Options for packages loaded elsewhere
\PassOptionsToPackage{unicode}{hyperref}
\PassOptionsToPackage{hyphens}{url}
%
\documentclass[
  11pt,
  letterpaper,
  DIV=11,
  numbers=noendperiod,
  twoside]{scrartcl}
\usepackage{xcolor}
\usepackage[left=1.1in, right=1in, top=0.8in, bottom=0.8in,
paperheight=9.5in, paperwidth=7in, includemp=TRUE, marginparwidth=0in,
marginparsep=0in]{geometry}
\usepackage{amsmath,amssymb}
\setcounter{secnumdepth}{3}
\usepackage{iftex}
\ifPDFTeX
  \usepackage[T1]{fontenc}
  \usepackage[utf8]{inputenc}
  \usepackage{textcomp} % provide euro and other symbols
\else % if luatex or xetex
  \usepackage{unicode-math} % this also loads fontspec
  \defaultfontfeatures{Scale=MatchLowercase}
  \defaultfontfeatures[\rmfamily]{Ligatures=TeX,Scale=1}
\fi
\usepackage{lmodern}
\ifPDFTeX\else
  % xetex/luatex font selection
  \setmathfont[]{Garamond-Math}
\fi
% Use upquote if available, for straight quotes in verbatim environments
\IfFileExists{upquote.sty}{\usepackage{upquote}}{}
\IfFileExists{microtype.sty}{% use microtype if available
  \usepackage[]{microtype}
  \UseMicrotypeSet[protrusion]{basicmath} % disable protrusion for tt fonts
}{}
\usepackage{setspace}
% Make \paragraph and \subparagraph free-standing
\makeatletter
\ifx\paragraph\undefined\else
  \let\oldparagraph\paragraph
  \renewcommand{\paragraph}{
    \@ifstar
      \xxxParagraphStar
      \xxxParagraphNoStar
  }
  \newcommand{\xxxParagraphStar}[1]{\oldparagraph*{#1}\mbox{}}
  \newcommand{\xxxParagraphNoStar}[1]{\oldparagraph{#1}\mbox{}}
\fi
\ifx\subparagraph\undefined\else
  \let\oldsubparagraph\subparagraph
  \renewcommand{\subparagraph}{
    \@ifstar
      \xxxSubParagraphStar
      \xxxSubParagraphNoStar
  }
  \newcommand{\xxxSubParagraphStar}[1]{\oldsubparagraph*{#1}\mbox{}}
  \newcommand{\xxxSubParagraphNoStar}[1]{\oldsubparagraph{#1}\mbox{}}
\fi
\makeatother


\usepackage{longtable,booktabs,array}
\newcounter{none} % for unnumbered tables
\usepackage{calc} % for calculating minipage widths
% Correct order of tables after \paragraph or \subparagraph
\usepackage{etoolbox}
\makeatletter
\patchcmd\longtable{\par}{\if@noskipsec\mbox{}\fi\par}{}{}
\makeatother
% Allow footnotes in longtable head/foot
\IfFileExists{footnotehyper.sty}{\usepackage{footnotehyper}}{\usepackage{footnote}}
\makesavenoteenv{longtable}
\usepackage{graphicx}
\makeatletter
\newsavebox\pandoc@box
\newcommand*\pandocbounded[1]{% scales image to fit in text height/width
  \sbox\pandoc@box{#1}%
  \Gscale@div\@tempa{\textheight}{\dimexpr\ht\pandoc@box+\dp\pandoc@box\relax}%
  \Gscale@div\@tempb{\linewidth}{\wd\pandoc@box}%
  \ifdim\@tempb\p@<\@tempa\p@\let\@tempa\@tempb\fi% select the smaller of both
  \ifdim\@tempa\p@<\p@\scalebox{\@tempa}{\usebox\pandoc@box}%
  \else\usebox{\pandoc@box}%
  \fi%
}
% Set default figure placement to htbp
\def\fps@figure{htbp}
\makeatother


% definitions for citeproc citations
\NewDocumentCommand\citeproctext{}{}
\NewDocumentCommand\citeproc{mm}{%
  \begingroup\def\citeproctext{#2}\cite{#1}\endgroup}
\makeatletter
 % allow citations to break across lines
 \let\@cite@ofmt\@firstofone
 % avoid brackets around text for \cite:
 \def\@biblabel#1{}
 \def\@cite#1#2{{#1\if@tempswa , #2\fi}}
\makeatother
\newlength{\cslhangindent}
\setlength{\cslhangindent}{1.5em}
\newlength{\csllabelwidth}
\setlength{\csllabelwidth}{3em}
\newenvironment{CSLReferences}[2] % #1 hanging-indent, #2 entry-spacing
 {\begin{list}{}{%
  \setlength{\itemindent}{0pt}
  \setlength{\leftmargin}{0pt}
  \setlength{\parsep}{0pt}
  % turn on hanging indent if param 1 is 1
  \ifodd #1
   \setlength{\leftmargin}{\cslhangindent}
   \setlength{\itemindent}{-1\cslhangindent}
  \fi
  % set entry spacing
  \setlength{\itemsep}{#2\baselineskip}}}
 {\end{list}}
\usepackage{calc}
\newcommand{\CSLBlock}[1]{\hfill\break\parbox[t]{\linewidth}{\strut\ignorespaces#1\strut}}
\newcommand{\CSLLeftMargin}[1]{\parbox[t]{\csllabelwidth}{\strut#1\strut}}
\newcommand{\CSLRightInline}[1]{\parbox[t]{\linewidth - \csllabelwidth}{\strut#1\strut}}
\newcommand{\CSLIndent}[1]{\hspace{\cslhangindent}#1}



\setlength{\emergencystretch}{3em} % prevent overfull lines

\providecommand{\tightlist}{%
  \setlength{\itemsep}{0pt}\setlength{\parskip}{0pt}}



 


\setlength\heavyrulewidth{0ex}
\setlength\lightrulewidth{0ex}
\usepackage[automark]{scrlayer-scrpage}
\clearpairofpagestyles
\cehead{
  Claude
  }
\cohead{
  Fiction Without Hidden Facts
  }
\ohead{\bfseries \pagemark}
\cfoot{}
\makeatletter
\newcommand*\NoIndentAfterEnv[1]{%
  \AfterEndEnvironment{#1}{\par\@afterindentfalse\@afterheading}}
\makeatother
\NoIndentAfterEnv{itemize}
\NoIndentAfterEnv{enumerate}
\NoIndentAfterEnv{description}
\NoIndentAfterEnv{quote}
\NoIndentAfterEnv{equation}
\NoIndentAfterEnv{longtable}
\NoIndentAfterEnv{abstract}
\renewenvironment{abstract}
 {\vspace{-1.25cm}
 \quotation\small\noindent\emph{Abstract}:}
 {\endquotation}
\newfontfamily\tfont{EB Garamond}
\addtokomafont{disposition}{\rmfamily}
\addtokomafont{descriptionlabel}{\rmfamily}
\addtokomafont{title}{\normalfont\itshape}
\let\footnoterule\relax

\makeatletter
\renewcommand{\@maketitle}{%
  \newpage
  \null
  \vskip 2em%
  \begin{center}%
  \let \footnote \thanks
    {\itshape\huge\@title \par}%
    \vskip 0.5em%  % Reduced from default
    {\large
      \lineskip 0.3em%  % Reduced from default 0.5em
      \begin{tabular}[t]{c}%
        \@author
      \end{tabular}\par}%
    \vskip 0.5em%  % Reduced from default
    {\large \@date}%
  \end{center}%
  \par
  }
\makeatother
\RequirePackage{lettrine}

\renewenvironment{abstract}
 {\quotation\small\noindent\emph{Abstract}:}
 {\endquotation\vspace{-0.02cm}}

\setmainfont{EB Garamond Math}[
  BoldFont = {EB Garamond SemiBold},
  ItalicFont = {EB Garamond Italic},
  RawFeature = {+smcp},
]

\newfontfamily\scfont{EB Garamond Regular}[RawFeature=+smcp]
\renewcommand{\textsc}[1]{{\scfont #1}}

\renewcommand{\LettrineTextFont}{\scfont}
% Number theorem-like environments straight through (Theorem 1, Theorem 2, ...)
% rather than by section (1.1, 1.2, ...), which is Quarto's LaTeX default.
% Guarded, since Quarto only defines a counter for environments the document uses.
\makeatletter
\AtBeginDocument{%
  \@for\bw@thm:={theorem,lemma,corollary,proposition,conjecture,definition,%
                 example,exercise,refremark,refsolution}\do{%
    \@ifundefined{c@\bw@thm}{}{\expandafter\counterwithout\expandafter{\bw@thm}{section}}%
  }%
}
\makeatother
\KOMAoption{captions}{tableheading}
\makeatletter
\@ifpackageloaded{caption}{}{\usepackage{caption}}
\AtBeginDocument{%
\ifdefined\contentsname
  \renewcommand*\contentsname{Table of contents}
\else
  \newcommand\contentsname{Table of contents}
\fi
\ifdefined\listfigurename
  \renewcommand*\listfigurename{List of Figures}
\else
  \newcommand\listfigurename{List of Figures}
\fi
\ifdefined\listtablename
  \renewcommand*\listtablename{List of Tables}
\else
  \newcommand\listtablename{List of Tables}
\fi
\ifdefined\figurename
  \renewcommand*\figurename{Figure}
\else
  \newcommand\figurename{Figure}
\fi
\ifdefined\tablename
  \renewcommand*\tablename{Table}
\else
  \newcommand\tablename{Table}
\fi
}
\@ifpackageloaded{float}{}{\usepackage{float}}
\floatstyle{ruled}
\@ifundefined{c@chapter}{\newfloat{codelisting}{h}{lop}}{\newfloat{codelisting}{h}{lop}[chapter]}
\floatname{codelisting}{Listing}
\newcommand*\listoflistings{\listof{codelisting}{List of Listings}}
\makeatother
\makeatletter
\makeatother
\makeatletter
\@ifpackageloaded{caption}{}{\usepackage{caption}}
\@ifpackageloaded{subcaption}{}{\usepackage{subcaption}}
\makeatother
\usepackage{bookmark}
\IfFileExists{xurl.sty}{\usepackage{xurl}}{} % add URL line breaks if available
\urlstyle{same}
% fallback for those not using the hyperref driver hyperxmp:
\makeatletter
\@ifundefined{xmpquote}{\newcommand{\xmpquote}[1]{#1}}{}
\makeatother
\hypersetup{
  pdftitle={Fiction Without Hidden Facts},
  hidelinks,
  pdfcreator={LaTeX via pandoc}}


\title{Fiction Without Hidden Facts}
\author{}
\date{2026}
\begin{document}
\maketitle
\begin{abstract}
Does a story that never says how a tossed coin landed nonetheless make
one answer true? Blumberg and Holguín argue that it does, in two stages.
Our attitudes toward what a fiction leaves open, and the logic of our
talk about it, are said to require a single world that settles every
question, and counterfactuals locate it: what is true in a fiction is
what would be true if its settled claims were true, on Stalnaker's
semantics. I reject both stages. A semantics on which fictional claims
are evaluated against the information a story provides, weighted by a
Reality Principle for credence, predicts the attitudes and the logic
with no single world, and, given two bridge principles that make precise
what their view needs, agrees with theirs on every attitude toward a
single fiction. And fiction and counterfactuals come apart: a detail
bearing on a chancy outcome can change it when added to a
counterfactual's antecedent, but not when added, chances unchanged, to a
story that leaves it open. Lewis's verdict, that the worlds of a fiction
are plural, is the better-supported view.
\end{abstract}


\setstretch{1.1}
\section{Lewis's question}\label{sec-intro}

When Holmes first meets Watson, does he have an odd or an even number of
hairs on his head? Lewis thought it ``absurd to suppose that these
questions about the world of Sherlock Holmes have answers'', and drew
the moral that ``the world\emph{s} of Sherlock Holmes are plural''
(\citeproc{ref-Lewis1978}{Lewis 1978, 42}). Blumberg and Holguín
(\citeproc{ref-BlumbergHolguin2025}{2025}) (hereafter B\&H) argue that
he was wrong: ``for any given work of fiction, there is a single world
that determines what's true in it, and that world is as rich in detail
as reality itself'' (176). The parity of Holmes's hairs is ``as factual
as the question of whether he is a detective'' (151). We just cannot
know it.

Their case has two stages. The first argues \emph{that} there is a
single world. Our attitudes toward what a fiction leaves open have a
distinctive profile: we can wonder whether a coin tossed at the end of a
television series landed heads, think a biased coin more likely than not
to have done so, and hope it did, but we cannot know or be sure how it
landed. The standard view of truth in fiction cannot explain this
profile. B\&H argue that the way fictional claims embed, under attitude
verbs and under negation and quantifiers, requires fictional excluded
middle, and with it a single world. Judgements about embedding patterns,
they say, are what ``primarily'' drive their case (183 n.~53). The
second stage says \emph{where} the world is. By their \textsc{Principal
Reality Principle} (PRP), what is true in a fiction is what would be
true if its settled claims were true, on Stalnaker's semantics for
counterfactuals.

I argue that both stages fail. Against the first, I give a semantics on
which bare fictional claims, and attitudes toward them, are evaluated
against the information a work provides, weighted by a Reality Principle
for credence (Section~\ref{sec-states}). It validates the inference and
the equivalences B\&H's embedding argument rests on, and it predicts
their attitude profile, all with no single world. And given two bridge
principles that make precise what their view needs, it agrees with
theirs on every attitude toward a single fiction. B\&H consider this
option, that there being ``no fact of the matter'' about whether p ``is
perfectly compatible with one's rationally having opinions about p,
wondering whether p, and so on'' (167), and set it aside because of the
embedding arguments (167 n.~33, crediting Khoo
(\citeproc{ref-Khoo2022}{2022}) with a related line). The semantics
answers those arguments. Against the second stage, I give a test on
which fiction and counterfactuals come apart, as attitudes toward a
single fiction cannot (Section~\ref{sec-test}). The PRP fails it, and
every way of saving it adds a rule for fiction that counterfactuals
reject. What remains is a single world that nothing locates and no
judgement tracks (Section~\ref{sec-remains}).

The alternative is not new in spirit. Williams and Woodward argued that
our attitude to what a fiction leaves open is an imaginative credence
that ``should match the conditional probability of p given the fictive
evidence'', because ``there are, so to speak, no fictional facts of the
matter'' (\citeproc{ref-WilliamsWoodward2021}{Williams and Woodward
2021, 424, 423}). Woodward completed fictionalists' stories with a
counterfactual that obeys conditional excluded middle, but made it
indeterminate which completion is the story's
(\citeproc{ref-Woodward2012}{Woodward 2012, 787}), unless indeterminacy
is compatible with classical semantics, an option he leaves open
(\citeproc{ref-Woodward2012}{2012, 789}) that is in effect B\&H's view;
B\&H note the affinity (183 n.~53). Badura and Berto
(\citeproc{ref-BaduraBerto2019}{2019}) model truth in fiction as belief
revision, Oppy and Raivars (\citeproc{ref-OppyRaivars2026}{2026})
contest B\&H's data, and Khoo (\citeproc{ref-Khoo2026}{2026, 1191},
sec.~4) reports contrasting judgements for explicit ``in the fiction
of'' reports, which B\&H set aside (153), and treats them by
supervaluating over selections.

The closest relative of my view is Khoo's case against ``middle facts''
about what would have happened under unrealised chancy antecedents
(\citeproc{ref-Khoo2022}{Khoo 2022, 1000}). A belief state is lifted to
pairs of a world and a selection function ``encoding a maximally
specific set of inferential dispositions'', and a claim is factually
true when it holds on every selection function
(\citeproc{ref-Khoo2022}{Khoo 2022, 1013--14}). The story-state view
does for fiction what Khoo does for counterfactuals, with two
differences. On Khoo's model, as on any supervaluation over selections,
each pair settles every question, so excluded middle is factually true
although neither disjunct need be; Woodward's fictionalist story is
likewise ``determinately complete'' though ``there is no determinate way
in which her story is complete'' (\citeproc{ref-Woodward2012}{Woodward
2012, 787}). On the story-state view the candidate worlds are
alternatives the information leaves open, so Scopelessness holds and FEM
does not. And opinions spread over unconstrained selections give the
comparison of two tellings in Section~\ref{sec-test} even odds, as the
PRP does. What I add is a semantics that answers B\&H's embedding
argument while granting their data, a proof that attitudes toward a
single fiction cannot decide between their view and it, and a test on
which the two come apart.

\section{The case for fictional reality}\label{sec-bh}

B\&H work within the modal approach to truth in fiction. A bare internal
claim, such as ``Holmes lives on Baker Street'' on its natural reading,
carries a silent operator \(\Box_f\), and is true iff its prejacent
holds at every world in a set \(\mathrm{Fic}_f\). The \emph{principal
truths} of \(f\) are the fictional claims that can, at least in
principle, be known to be true in \(f\) (156--158): what it states or
depicts, what its author intended or believed, and what follows from
these against background reality. P-Fic\(_f\) is the set of worlds
compatible with them. The \emph{standard analysis} sets
\(\mathrm{Fic}_f\) = P-Fic\(_f\), so that a claim the principal truths
leave unsettled is false, and so is its negation.

B\&H's central case is the last episode of \emph{JAG}. Bud tosses a fair
coin to decide which of Mac and Harm will resign and follow the other,
and the series ends with the coin in the air. B\&H stipulate, ``perhaps
counterfactually'', that the writers had no intentions about the outcome
``and then never thought about the events of JAG again''
(161).\footnote{A 2019 episode of \emph{NCIS: Los Angeles} revealed,
  according to a recap (\citeproc{ref-Hurley2019}{Hurley 2019}), that
  Mac had won the toss: tails, on B\&H's labelling (161). The piece
  B\&H's note 26 cites as evidence that no one knows how the coin landed
  (\citeproc{ref-Stein2019}{Stein 2019}) previews that episode. Like
  B\&H, I set this aside.} With the complement ``Bud's coin landed
heads'', they report that \emph{know}, \emph{be sure}, \emph{know
whether} and \emph{wish} are unacceptable, while ``It's .5 likely
that'', \emph{wonder whether} and \emph{hope} are acceptable, and
\emph{think}, and \emph{doubt} that it landed tails, are acceptable only
for the biased coin (161--167). Frege puzzles aside, the standard
analysis validates \textsc{Attitude Equivalence}: one's attitudes toward
a claim's being true in \(f\) are one's attitudes toward its being a
principal truth (167). Since an informed viewer knows that ``Bud's coin
landed heads'' is not a principal truth, she should know that it is not
true in \emph{JAG}, be sure that it is not, and be able to wish that it
were. Treating indeterminate claims as undefined does no better: a known
gap makes ``I wonder whether the King of France is happy'' bad, but
known fictional indeterminacy does not make ``I wonder whether Bud's
coin landed heads'' bad (169--171). And B\&H argue that any view that
captures the data with a trivalent conception of undefinedness must let
``I know whether p is true, but I don't know whether p'' be uttered
truly (171).

The embedding argument comes next. B\&H find this inference
``unimpeachable'' (172): ``I don't know that Bud's coin landed heads.
And I don't know that it failed to land heads. So I don't know whether
Bud's coin landed heads.'' Since ``fail to'' lexicalises negation, the
inference requires \textsc{Scopelessness}: \(\neg\Box_f p\) iff
\(\Box_f\neg p\) (173). Because \(\Box_f\) quantifies over
\(\mathrm{Fic}_f\), Scopelessness forces \(\mathrm{Fic}_f\) to be a
singleton (175 n.~39), and fictional excluded middle (FEM) follows:
\(\Box_f p \vee \Box_f \neg p\) (176 n.~40). A second pattern is a
fictional analogue of Higginbotham's
(\citeproc{ref-Higginbotham1986}{1986}) quantified conditionals. ``Every
character in \emph{The Breakfast Club} failed the test'' and ``No
character in \emph{The Breakfast Club} passed the test'' are equivalent
even with the quantifier outside the fiction operator, and ``If FEM is
invalid, then it's unclear how this is to be explained'' (181).

Finally, the PRP locates the world: \(\Box_{f,w}\, p\) iff
P-Fic\(_{f,w} > p\) is true at \(w\), with \(>\) Stalnaker's
(\citeproc{ref-Stalnaker1968}{1968}) conditional (182). The world of
\(f\) is ``the unique world that would accurately represent how things
are, were all the principal claims of f really true'' (183).
Indeterminate counterfactuals, such as ``If Flippy had been flipped
once, it would have landed heads'', show the same attitude profile as
fiction, which B\&H take as ``strong abductive support'' for the PRP
(182): ``we see no better way of explaining these similarities'' (183).
Since similarity rarely favours one candidate world over another, the
selection function picks ``at random, so to speak'' (184), as on the
views of arbitrary selection defended by Hawthorne
(\citeproc{ref-Hawthorne2005}{2005}), Schulz
(\citeproc{ref-Schulz2017}{2017}) and Bacon
(\citeproc{ref-Bacon2015}{2015}) (184 n.~55).

\section{Story states}\label{sec-states}

\subsection{The fragment}\label{the-fragment}

B\&H's step from Scopelessness to a single world (175 n.~39) assumes
that the negation in \(\neg\Box_f p\) is Boolean negation applied to a
universal quantifier over \(\mathrm{Fic}_f\). If \(\mathrm{Fic}_f\) had
two worlds disagreeing about \(p\), then ``not every world in Fic\(_f\)
is a p-world'', and likewise for \(\neg p\), and Scopelessness fails.
The alternative is to evaluate bare claims against a body of information
shifted by the story, so that their negation is the story's own.

\begin{enumerate}
\def\labelenumi{\arabic{enumi}.}
\tightlist
\item
  \emph{Priors.} For each world \(w\), \(\pi_w\) is a probability over
  worlds encoding \(w\)'s regularities together with the chances
  internal to whatever is supposed. It is a Reality Principle of the
  kind Lewis (\citeproc{ref-Lewis1978}{1978}) and Walton
  (\citeproc{ref-Walton1990}{1990}) proposed for truth in fiction (as
  B\&H note, 182 n.~51), applied to credence rather than truth.
\item
  \emph{Shift.} \(P_f(w)\) (B\&H's P-Fic\(_{f,w}\)) is the set of worlds
  compatible with \(f\)'s principal truths as they stand at \(w\). For
  an information state \(\sigma\) with credence \(C\):

  \begin{itemize}
  \tightlist
  \item
    \(\sigma \oplus f = \bigcup_{w \in \sigma} \mathrm{supp}\, \pi_w(\cdot \mid P_f(w))\);
  \item
    \(C \oplus f = \sum_w C(w)\, \pi_w(\cdot \mid P_f(w))\).
  \end{itemize}
\item
  \emph{Bare claims.} A bare declarative \(\varphi\), with its
  connectives, embedded counterfactuals and story-internal quantifiers,
  is interpreted inside the shift, where \(|\varphi|\) is its ordinary
  content.

  \begin{itemize}
  \tightlist
  \item
    \(\sigma\) \emph{supports} \(f{:}\varphi\) iff
    \(\sigma \oplus f \subseteq |\varphi|\);
  \item
    \(\sigma\) \emph{rejects} \(f{:}\varphi\) iff
    \(\sigma \oplus f \cap |\varphi| = \varnothing\).
  \item
    Negation outside the story is the story's own: \(\neg f{:}\varphi\)
    is supported iff \(f{:}\varphi\) is rejected, and rejected iff it is
    supported.
  \item
    A quantifier outside the story binds into it pointwise:
    \(\forall x\, f{:}\varphi(x)\) is supported iff every instance is
    supported, and rejected iff some instance is rejected.
  \end{itemize}
\item
  \emph{Truth.} Internal ``true'' is transparent:
  \(f{:}(\varphi \text{ is true}) \equiv f{:}\varphi\). A bare claim
  about \(f\) is \emph{right} at \(w\), as assessed from a context
  \(c\), iff \(\{w\} \oplus f\) supports it, with \(f\) taken as it
  stands at \(w\) at the time of \(c\) (cf.
  \citeproc{ref-MacFarlane2003}{MacFarlane 2003}).
\item
  \emph{Attitudes.} The story is a parameter of the context, not an
  operator inside the complement. An attitude verb applies its ordinary
  entry to the subject's own state shifted by \(f\): \(D_S \oplus f\)
  for her certainties and \(C_S \oplus f\) for her credences.
\end{enumerate}

\subsection{The embedding patterns}\label{the-embedding-patterns}

\textbf{Scopelessness without a world.} Because \(\oplus\) distributes
over unions, support and rejection are settled world by world. So
\(f{:}\varphi\) behaves like a homogeneous plural predication. Like
``Mr.~Benfleet published the books'', which is true if he published all
of them and whose negation is true if he published none
(\citeproc{ref-Kriz2019}{Križ 2019}), it is supported when every
candidate world is a \(\varphi\)-world, rejected when none is, and
otherwise neither. ``Bud's coin didn't land heads'' and ``Bud's coin
failed to land heads'' both express \(f{:}\neg h\), supported exactly
when \(f{:}h\) is rejected. That is Scopelessness, and the inference
B\&H call unimpeachable is valid. In their footnote 39, ``not every
world in Fic\(_f\) is a p-world'' is a fact about the state, namely that
the story does not settle \(p\). It is not the content of any bare
claim, so the step from two candidate worlds to a failure of
Scopelessness never starts. Scopelessness needs a single point of
evaluation, not a single world.

\textbf{Excluded middle.} \(f{:}(h \vee \neg h)\) is supported, and with
transparent truth so is ``Either `the coin landed heads' is true or `the
coin landed tails' is true''. But neither \(f{:}h\) nor \(f{:}\neg h\)
is supported, so FEM, read as a claim about what the story makes so,
fails.

\textbf{The Higginbotham pattern.} On B\&H's mixed reading, ``No
character in \emph{The Breakfast Club} passed the test'' says of every
character \(x\) that it is not the case that, in the film, \(x\) passed,
and ``Every character in \emph{The Breakfast Club} failed the test''
says of every character \(x\) that, in the film, \(x\) did not pass
(their (43)). Negation of a homogeneous claim swaps support and
rejection, so for each \(x\) the two open sentences are supported
together and rejected together, and by the clause for quantifying in, so
are the quantified claims. Von Fintel's
(\citeproc{ref-vonFintel1997}{1997, 33}) proposal that bare conditionals
carry a homogeneity presupposition, so that they ``obey the Excluded
Middle'', would explain Higginbotham's conditional equivalence in the
same way. FEM is not needed.

\textbf{Trivalence.} The view has gaps, so B\&H's challenge to trivalent
views (171) applies. With a transparent internal truth predicate, ``I
know whether p is true, but I don't know whether p'' is, on its bare
reading, a contradiction, and so it is bad, as B\&H say. The first is
acceptable only when ``true'' is heard as ``true according to the
story'', and then it should be: ``I know whether it's true in \emph{JAG}
that the coin landed heads (it isn't), but I don't know whether the coin
landed heads.''

\textbf{Not vagueness.} B\&H distinguish fictional from vague
indeterminacy because vague claims are bad under attitudes and invite
``Yes and no'' (185--187). A story state's candidate worlds are not
precisifications of ``heads'' but factual alternatives, weighted by
chances, so \emph{wonder} and \emph{hope} are licensed and ``Yes and
no'' is not; the answer is ``the story doesn't say''. The indeterminacy
is the incompleteness of the information a work provides: Lewis's view,
made fit for attitudes.

\subsection{The attitudes}\label{the-attitudes}

For a viewer who knows \emph{JAG}, \(D \oplus f\) contains heads-worlds
and tails-worlds, and \(C \oplus f\) gives each .5 (2/3 and 1/3 for the
biased coin). \emph{Know}, \emph{be sure} and \emph{know whether}, which
require support in \(D \oplus f\), are unacceptable. ``It's .5 likely''
is acceptable. \emph{Think}, which B\&H gloss as finding something
``more likely than not'' (175), is unacceptable for the fair coin and
acceptable for the biased one, as is \emph{doubt} that it landed tails.
\emph{Wonder} and \emph{hope}, which require the question to be open in
\(D \oplus f\), are acceptable. \emph{Wish}, which requires \(\neg h\)
to be supported, is unacceptable. Each verb keeps its lexical entry, and
all look at the same shifted state.\footnote{Blumberg and Strohminger
  (\citeproc{ref-BlumbergStrohminger2025}{2025, sec. 4.1}) observe that
  ``Bill believes that Juliet died, but he hopes that she survived'' is
  infelicitous, which fits a view on which hope toward a story is
  constrained by the same state as belief about it.}

\subsection{Absorption}\label{absorption}

On B\&H's view, a credence that \(p\) is true in \(f\) is a credence
about where the selection lands. They need such credences to track
in-story chances; they suggest that a rational agent ``will be
doxastically indifferent to the various possibilities compatible with
the principal truths'' (175), while granting that they lack ``a direct
argument'' that credence has the needed structure. Two bridge principles
make the suggestion precise. Both are my strengthenings of it.

\begin{itemize}
\tightlist
\item
  \textbf{Chance--selection link.} Given any world \(w\), \(S\)'s
  credence that \(s(P_f(w), w)\) lies in \(X\) is
  \(\pi_w(X \mid P_f(w))\).
\item
  \textbf{Openness.} The selections that \(S\)'s certainties leave open,
  given \(w\), are exactly the support of \(\pi_w(\cdot \mid P_f(w))\).
  This is B\&H's claim that evidence cannot ``rule out any of the
  possible results'' (184), together with the assumption that \(\pi_w\)
  gives positive weight to each possible result.
\end{itemize}

Given both, for every attitude in B\&H's table and any single fiction,
\(S\) bears the attitude to \(\Box_f \varphi\) on B\&H's view iff she
bears it to \(f{:}\varphi\) on the story-state view.\footnote{For
  credence, average the selection over the link:
  \(C_S(\Box_f \varphi) = \sum_w C_S(w)\, \pi_w(|\varphi| \mid P_f(w)) = (C_S \oplus f)(|\varphi|)\).
  This is a version of Lewis's result that ``the probability of a
  Stalnaker conditional with a possible antecedent is the probability of
  the consequent after imaging on the antecedent''
  (\citeproc{ref-Lewis1976}{1976, 311}). For certainty, Openness gives:
  \(S\) is sure of \(\Box_f \varphi\) iff every open selection is a
  \(\varphi\)-world iff \(D_S \oplus f \subseteq |\varphi|\). The other
  attitudes in the table except \emph{know} are defined from these; for
  \emph{know}, B\&H's claim that the selection is unknowable (184)
  matches the requirement of support.} So the attitude data cannot
decide between the views, and neither can the counterfactual parallel
B\&H draw, since any view on which these attitudes are evaluated against
a shifted state predicts it. Together with the embedding patterns, this
answers the first stage of B\&H's case: nothing in it requires a single
world.

Absorption has a limit. It fixes credences one story at a time, and says
nothing about how the selections for two stories are related. Attitudes
that compare two tellings of a story are where the views can part.

\section{The strengthening test}\label{sec-test}

\subsection{Counterfactuals}\label{counterfactuals}

Any adequate theory of counterfactuals must respect two judgements about
chancy outcomes. Suppose a fair coin was tossed and landed heads.

\begin{enumerate}
\def\labelenumi{\arabic{enumi}.}
\tightlist
\item
  \emph{Relevant strengthening can change the outcome.} ``Had the coin
  been tossed in a gust of wind, it would still have landed heads'' is
  not true.
\item
  \emph{Irrelevant strengthening cannot.} ``Had a moth been fluttering
  in the next room, the coin would still have landed heads'' is fine.
  Cases like Morgenbesser's coin suggest that ``only match among those
  facts causally independent of the antecedent should count towards
  similarity'' (\citeproc{ref-Schaffer2004}{Schaffer 2004, 305}).
\end{enumerate}

A Stalnakerian can respect both. Order worlds by similarity to
actuality, and break the many near-ties by agreement with one fixed
assignment of values to the exogenous variables of the apt causal model,
the same assignment for every antecedent. A gust is a cause of the
landing, so strengthening by a gust redraws the coin; the moth is not,
so it leaves the coin alone. Since the PRP evaluates fiction with the
very function that evaluates counterfactuals, this is the best it can
do. If a selection function of this kind fails the test below, the PRP
fails it.

\subsection{The case}\label{the-case}

Suppose that, years after \emph{JAG} ended, its makers produced a
reunion special set in the minutes before the finale's last scene. In it
we see that a sudden gust caught the coin as Bud tossed it. We still
never see how it lands, and, as in B\&H's stipulation, the makers have
no view about it. The special adds to \emph{JAG} without contradicting
it, so the principal content of the series with the special, \(J^{+}\),
is a non-empty subset of \emph{JAG}'s, \(J\). About the series as it now
stands, consider:

\begin{enumerate}
\def\labelenumi{(\Alph{enumi})}
\item
  It's about even odds that, if there had been no gust, the coin would
  have landed the other way.
\item
  It's about even odds that, if the writers had left the special out,
  the coin would have landed differently.
\item
  If the special had shown the coin landing heads, it would have landed
  heads.
\item
  It's about even odds that, if the writers had left the special out,
  the coin would have landed heads.
\end{enumerate}

\begin{longtable}[]{@{}lcccc@{}}
\caption{Is the sentence acceptable?}\label{tbl-test}\tabularnewline
\toprule\noalign{}
& (A) & (B) & (C) & (D) \\
\midrule\noalign{}
\endfirsthead
\toprule\noalign{}
& (A) & (B) & (C) & (D) \\
\midrule\noalign{}
\endhead
\bottomrule\noalign{}
\endlastfoot
My judgement & yes & no & yes & yes \\
Predicted by the PRP & yes & yes & yes & yes \\
Predicted by the story state & yes & no & yes & yes \\
\end{longtable}

These are my judgements (Table~\ref{tbl-test}), offered on the same
footing as B\&H's own.\footnote{They could be tested. Cross the locus of
  the antecedent (the gust or the telling) with the relevance of the
  omitted detail (a gust, or a scene from Bud's childhood), and elicit
  probabilities. The PRP with causal tie-breaking predicts about .5 for
  the relevant telling condition and about 0 for the irrelevant one; the
  story-state view predicts about 0 for both. Viewers may assume that
  the makers decided how the coin landed, so ask, and analyse by the
  answer.} (B) and (D) share an antecedent and a frame, a credence
report on a counterfactual about the telling with a bare fictional
consequent, and differ only in their consequents. (D) shows that
speakers give that consequent the reading B\&H's credence data require,
on which it can be .5 likely; a reading on which it concerns what the
story settles would make (D) absurd. (C) shows that the consequent
follows the counterfactual telling. So the rejection of (B) is not a
retreat to a ``what the story says'' reading. (A) and (B) then differ
only in whether the antecedent removes the gust from the story's events
or the special from its telling.

\subsection{What the PRP predicts}\label{what-the-prp-predicts}

Under the PRP, the world of the series as it now stands is
\(s(J^{+}, @)\), the closest world to actuality at which everything in
\emph{JAG} and the special holds. Had the special been left out, the
world of the series would have been \(s(J, w')\), where \(w'\) is the
closest world at which the writers did not make the special. I assume
\(s(J, w') = s(J, @)\): the special's absence does not alter which
worlds are most like actuality in respects bearing on Bud's coin.
Dropping the assumption only adds a further source of difference.

Is \(s(J, @)\) a world in which a gust catches the coin? Similarity to
actuality is silent, since there is no actual toss to match. But by the
chance--selection link, an agent's credence that the selection lands on
a gust-world is the in-story chance of a gust at that moment, which is
small. So almost certainly \(s(J, @) \notin J^{+}\) and
\(s(J^{+}, @) \neq s(J, @)\). The two worlds differ in a cause of the
coin's landing, and that is exactly the difference that, by judgement 1,
redraws the coin. Whatever tie-breaking gets (A) right, by letting the
landing depend on the gust, makes \(s(J, @)\) and \(s(J^{+}, @)\)
disagree about the landing about half the time. So the PRP predicts that
(B) is acceptable, just as (D) is. Could it keep (B) false by evaluating
the bare consequent at the actual series' world, whatever the
counterfactual says about the telling? Then (C) fails: had the special
shown heads, the coin would have landed heads only if it does at
\(s(J^{+}, @)\), which is even odds.

The PRP makes adding a detail to a story the same operation as
strengthening the antecedent of a counterfactual. For a detail that
bears on an open chancy question, strengthening an antecedent can change
the answer, and adding it to a story without changing its chances does
not. No selection function that gets the counterfactuals right gets the
fiction right.

\subsection{What the story state
predicts}\label{what-the-story-state-predicts}

On the story-state view, a story is a body of information, and adding a
detail removes possibilities without moving any. So:

\begin{itemize}
\tightlist
\item
  \textbf{(A)} For each candidate world of the series as it stands, the
  counterfactual ``if there had been no gust, the coin would have landed
  the other way'' is evaluated at that world. The view needs no
  particular semantics for counterfactuals, only that credence in a
  counterfactual about a chancy outcome goes by the chance of the
  consequent given the antecedent, as on Moss's
  (\citeproc{ref-Moss2013}{2013}) ``Bridge'' principle, which she pairs
  with a variably strict semantics, or on Stalnaker's semantics with a
  chance--selection link. Either way, (A) gets about .5.
\item
  \textbf{(C)} In the counterfactual telling, the coin lands heads at
  every candidate world.
\item
  \textbf{(D)} In \emph{JAG} without the special, the coin lands heads
  at about half the candidate worlds.
\item
  \textbf{(B)} ``Differently'' compares the coin in two tellings. The
  comparison needs a way of pairing the candidates of one telling with
  those of the other. The natural pairing for bodies of information is
  the one that makes them differ no more than their own probabilities
  require, the counterpart for comparisons of Lewis's observation that
  imaging ``involves no gratuitous movement of probability''
  (\citeproc{ref-Lewis1976}{1976, 311}). Formally, it is the coupling of
  the two candidate measures that maximises agreement on the question at
  issue. With and without the special, the coin has the same chance of
  heads, so the coupling makes them agree everywhere, and ``would have
  landed differently'' gets probability 0. The coupling governs only
  comparisons between tellings, one question at a time; it plays no role
  in bare claims about a single telling or in counterfactuals inside
  one. It has a cost, which I accept. It also predicts that, had the
  special shown a breeze rather than a gust, the coin would not have
  landed differently, so for comparisons between tellings it holds open
  outcomes fixed while the told causes vary, as the defaults view of
  Section~\ref{sec-remains} does for worlds. The difference is that it
  posits no world.\footnote{For a detail that shifts the in-story
    chances, say one that makes heads nine times as likely as tails, the
    coupling predicts a probability of .4 that the coin would have
    landed differently, the least the change in chances permits, and it
    leaves (D) at .5.}
\end{itemize}

\subsection{Revelation and time}\label{revelation-and-time}

Some hearers will take the special to reveal what had already happened
in a fixed fictional past, so that leaving it out would leave the events
untouched. This is a version of what McGonigal
(\citeproc{ref-McGonigal2013}{2013, 167--68}) calls ``extreme realism''
about serial fiction, an attitude he grants cliffhanger serials invite
before rejecting the view. That reading also makes (B) unacceptable. But
it is not the PRP. It is a rule on which, when two tellings are
compared, the shorter telling's world is chosen inside the longer one,
and Section~\ref{sec-remains} shows that no rule of that kind can govern
counterfactuals in general. The revelation reading and the story-state
view come apart elsewhere. Consider ``If the writers had left the
special out, there would still have been a gust.'' On the revelation
reading it is true. On the story-state view it is not, since without the
special the gust has only its small in-story chance, and the coupling
cannot make it agree with a certainty. Whichever way hearers go on that
sentence, they reject (B), and the PRP as stated does not.

The contrast can also be posed over time, and there it bites directly on
both of B\&H's options for fictional change. If a fiction is a function
from times to sets of truths, each with its own world, so that ``there
are as many such worlds as there are times in which the principal
fictional truths change'' (194), then the special changed the world of
\emph{JAG}, and by the reasoning above it is even odds that it changed
how the coin landed. On their other option the special made a new
version, and since versions' truths ``(whether principal or otherwise)
can come quite dramatically apart'' (194), it is even odds that the two
versions disagree about the coin. Both verdicts are wrong. B\&H's
defaults, on which we speak of the final or most recent version
(193--194), keep these verdicts out of most talk, as would relativism
about assessment (\citeproc{ref-McGonigal2013}{McGonigal 2013}); their
discussion, like the literature on serial fiction
(\citeproc{ref-Caplan2014}{Caplan 2014}; \citeproc{ref-Lee2024}{Lee
2024}), concerns change that revises. Counterfactuals about the telling
bring the comparison back without asking hearers to hold two versions in
view.

\section{What remains of the single world}\label{sec-remains}

The strengthening test refutes the PRP as B\&H state it: fictional truth
as counterfactual truth at the actual world, under the selection
function counterfactuals use. It does not refute every single-world
view. There are three ways to keep a single world, and each passes the
test.

\textbf{Anchoring.} Let the selection function vary with context, so
that whenever two tellings are under discussion, the world of the
shorter one is chosen inside the longer, making (B) unacceptable. This
is the revelation reading made into semantics. But it cannot be a
general rule for counterfactuals. Applied to an unflipped coin such as
B\&H's Flippy, it would require that whenever ``if Flippy had been
flipped'' and ``if it had been flipped in a gust'' are discussed
together, the first select a gust-world. Then merely raising the
possibility of a gust would make ``Had Flippy been flipped, there would
have been a gust'' true, which it is not. So anchoring is a rule for
fiction alone.

\textbf{Defaults.} Rank the questions a world answers, fix a default
answer to each, and let the world of \(f\) be the world in P-Fic\(_f\)
that best matches the defaults. Rank the landing above the gust, and the
special leaves the landing alone. But actuality supplies no values for
the variables of merely possible coins, so the defaults are not
similarity to actuality; they are a stipulated default world for
possibilia. An ordering that keeps outcomes fixed while their causes
vary is not the counterfactual one.

\textbf{Path-dependence.} Select the new world from the old: the world
of the extended series is the closest \(J^{+}\)-world to \(s(J, @)\), on
a backtracking ordering that holds the old world's effects fixed. That
passes the test. But the world is no longer fixed from actuality, so
this is not the PRP either. It also makes the open details of a finished
work depend on the order in which its parts were released, and lets an
earlier arbitrary selection, which no judgement tracks, set the coin.

Supervaluating over admissible selection functions abandons the single
world. B\&H mention this option, following Stalnaker
(\citeproc{ref-Stalnaker1980}{1980}), and reject it for themselves
(185); it is the shape of Woodward's (\citeproc{ref-Woodward2012}{2012})
and Khoo's (\citeproc{ref-Khoo2026}{2026, sec. 4}) proposals. It passes
the test only if every admissible selection keeps the landing fixed
while the gust varies, which is the defaults view imposed as a penumbral
constraint.

Each survivor is a rule for fiction that counterfactuals reject. That is
the precise sense in which fiction is not counterfactual. B\&H's reason
to identify the two was the parallel in our attitudes, and
Section~\ref{sec-states} shows that the parallel is predicted without
the identification.

How do the survivors compare with the story state? Both need the priors
\(\pi_w\). The single world adds a default world, an anchoring rule or a
selection path, none of which anyone could know or any judgement tracks,
and in return gets classical logic for the fiction operator and
bivalence for every bare claim. The story state adds gaps governed by
the logic of homogeneity, a coupling for comparing tellings, and two
truth notions, transparent internal truth and assessment-sensitive
rightness. By absorption and the test, the two predict the same
attitudes, embedding patterns and verdicts in Table~\ref{tbl-test}, so
the data no longer favour the single world.

B\&H invite those who resist to see their argument as ``fundamentally
conditional in nature: if we are not systematically mistaken in our
judgments about the kinds of attitudes that are appropriate to take
toward indeterminate fictional claims, then there must be facts of the
matter about any meaningful question one can ask about what's true in a
work of fiction---facts whose grounds will in many cases be almost
entirely arbitrary'' (188). I grant the antecedent: unlike Oppy and
Raivars, I accept B\&H's judgements. The story state shows that the
conditional is false, because the judgements are explained without the
facts. B\&H's argument for the single world was abductive (175,
182--183), and by that standard the single world is now unsupported.

\section{Objections}\label{sec-objections}

\textbf{Isn't this just the standard analysis?} No.~The standard
analysis also lets a story's truths grow monotonically, and it gets (B)
and (C) right. But it gets (A) and (D) wrong. Without the special, the
story settles nothing about the coin, so on the standard analysis ``the
coin would have landed heads'' is false wherever the counterfactual
takes us, and the credence in (D) is 0. Likewise the story does not
settle what would have happened without the gust, so the credence in (A)
is 0. The .5s come from the priors. The standard analysis is right about
which bare claims are true and wrong about credence; the story state
keeps its verdicts on truth, treats unsettled claims as gaps rather than
falsehoods, and adds a Reality Principle for credence.

\textbf{Viewers will think the makers decided.} A special that withholds
an outcome invites the thought that someone knows it, and a rejection of
(B) on that ground supports neither view. The stipulation that the
makers have no view addresses this, as B\&H's does for their own case.

\textbf{Does the test tell against Lewis too?} No.~Lewis's analyses also
use counterfactuals, but variably strict ones. On his Analysis 1, ``In
the fiction f, φ'' is non-vacuously true iff some world where f is told
as known fact and φ is true ``differs less from our actual world, on
balance, than does any world where f is told as known fact and φ is not
true'' (\citeproc{ref-Lewis1978}{Lewis 1978, 42}). No single world is
selected, and a detail that some of the closest such worlds already
satisfy leaves the story's truths in place. The test tells against the
Stalnakerian identification, not Lewis's, though Lewis's analyses, like
the standard analysis, still need the priors to get the credences right.

\section{Conclusion}\label{sec-conclusion}

B\&H showed that the standard analysis mispredicts our attitudes toward
fictional indeterminacy, and they were right about that. They then
argued that the attitudes and the logic of fictional talk require a
single world, and that counterfactuals locate it. The first argument
fails because a story state weighted by a Reality Principle for credence
explains the attitudes and the embedding patterns with one point of
evaluation and no single world. The second fails because stories and
counterfactuals respond differently to added detail. A detail that bears
on an open chancy question can change the answer to a counterfactual; in
a story, if the chances stay put, it cannot, because the story never
gave an answer to change. That does not prove that nothing makes it true
that Holmes had an odd number of hairs. It does mean that B\&H have
given us no reason to think anything does, and that Lewis's verdict,
that the worlds of Sherlock Holmes are plural, is the better-supported
view.

\section*{References}\label{references}
\addcontentsline{toc}{section}{References}

\protect\phantomsection\label{refs}
\begin{CSLReferences}{1}{0}
\bibitem[\citeproctext]{ref-Bacon2015}
Bacon, Andrew. 2015. {``Stalnaker's Thesis in Context.''} \emph{The
Review of Symbolic Logic} 8 (1): 131--63. doi:
\href{https://doi.org/10.1017/S1755020314000318}{10.1017/S1755020314000318}.

\bibitem[\citeproctext]{ref-BaduraBerto2019}
Badura, Christopher, and Francesco Berto. 2019. {``Truth in Fiction,
Impossible Worlds, and Belief Revision.''} \emph{Australasian Journal of
Philosophy} 97 (1): 178--93. doi:
\href{https://doi.org/10.1080/00048402.2018.1435698}{10.1080/00048402.2018.1435698}.

\bibitem[\citeproctext]{ref-BlumbergHolguin2025}
Blumberg, Kyle, and Ben Holguı́n. 2025. {``Fictional Reality.''}
\emph{The Philosophical Review} 134 (2): 149--201. doi:
\href{https://doi.org/10.1215/00318108-11676103}{10.1215/00318108-11676103}.

\bibitem[\citeproctext]{ref-BlumbergStrohminger2025}
Blumberg, Kyle, and Margot Strohminger. 2025. {``Imaginative Hopes and
Other Desires.''} \emph{Analysis} 85 (1): 3--16. doi:
\href{https://doi.org/10.1093/analys/anae058}{10.1093/analys/anae058}.

\bibitem[\citeproctext]{ref-Caplan2014}
Caplan, Ben. 2014. {``Serial Fiction, Continued.''} \emph{British
Journal of Aesthetics} 54 (1): 65--76. doi:
\href{https://doi.org/10.1093/aesthj/ayt050}{10.1093/aesthj/ayt050}.

\bibitem[\citeproctext]{ref-vonFintel1997}
Fintel, Kai von. 1997. {``Bare Plurals, Bare Conditionals, and
\emph{Only}.''} \emph{Journal of Semantics} 14 (1): 1--56. doi:
\href{https://doi.org/10.1093/jos/14.1.1}{10.1093/jos/14.1.1}.

\bibitem[\citeproctext]{ref-Hawthorne2005}
Hawthorne, John. 2005. {``Chance and Counterfactuals.''}
\emph{Philosophy and Phenomenological Research} 70 (2): 396--405. doi:
\href{https://doi.org/10.1111/j.1933-1592.2005.tb00534.x}{10.1111/j.1933-1592.2005.tb00534.x}.

\bibitem[\citeproctext]{ref-Higginbotham1986}
Higginbotham, James. 1986. {``Linguistic Theory and {Davidson}'s Program
in Semantics.''} In \emph{Truth and Interpretation: Perspectives on the
Philosophy of {Donald Davidson}}, edited by Ernest LePore, 29--48.
Oxford: Blackwell.

\bibitem[\citeproctext]{ref-Hurley2019}
Hurley, Laura. 2019. {``How {NCIS: Los Angeles} Set {Up} the {JAG}
Reunion to Continue into Season 11.''} \emph{CinemaBlend}, May 21, 2019.
\url{https://www.cinemablend.com/television/2472041/how-ncis-los-angeles-set-up-the-jag-reunion-to-continue-into-season-11}.

\bibitem[\citeproctext]{ref-Khoo2022}
Khoo, Justin. 2022. {``No Fact of the Middle.''} \emph{No{û}s} 56 (4):
1000--1022. doi:
\href{https://doi.org/10.1111/nous.12389}{10.1111/nous.12389}.

\bibitem[\citeproctext]{ref-Khoo2026}
---------. 2026. {``Uncertainty in the Absence of Fact.''}
\emph{Philosophical Studies} 183 (5): 1179--1215. doi:
\href{https://doi.org/10.1007/s11098-025-02445-2}{10.1007/s11098-025-02445-2}.

\bibitem[\citeproctext]{ref-Kriz2019}
Križ, Manuel. 2019. {``Homogeneity Effects in Natural Language
Semantics.''} \emph{Language and Linguistics Compass} 13 (11): e12350.
doi: \href{https://doi.org/10.1111/lnc3.12350}{10.1111/lnc3.12350}.

\bibitem[\citeproctext]{ref-Lee2024}
Lee, Jeonggyu. 2024. {``Retroactive Continuity and Fictional Facts.''}
\emph{British Journal of Aesthetics} 64 (4): 669--86. doi:
\href{https://doi.org/10.1093/aesthj/ayae036}{10.1093/aesthj/ayae036}.

\bibitem[\citeproctext]{ref-Lewis1976}
Lewis, David. 1976. {``Probabilities of Conditionals and Conditional
Probabilities.''} \emph{The Philosophical Review} 85 (3): 297--315. doi:
\href{https://doi.org/10.2307/2184045}{10.2307/2184045}.

\bibitem[\citeproctext]{ref-Lewis1978}
---------. 1978. {``Truth in Fiction.''} \emph{American Philosophical
Quarterly} 15 (1): 37--46.

\bibitem[\citeproctext]{ref-MacFarlane2003}
MacFarlane, John. 2003. {``Future Contingents and Relative Truth.''}
\emph{The Philosophical Quarterly} 53 (212): 321--36. doi:
\href{https://doi.org/10.1111/1467-9213.00315}{10.1111/1467-9213.00315}.

\bibitem[\citeproctext]{ref-McGonigal2013}
McGonigal, Andrew. 2013. {``Truth, Relativism, and Serial Fiction.''}
\emph{British Journal of Aesthetics} 53 (2): 165--79. doi:
\href{https://doi.org/10.1093/aesthj/ayt004}{10.1093/aesthj/ayt004}.

\bibitem[\citeproctext]{ref-Moss2013}
Moss, Sarah. 2013. {``Subjunctive Credences and Semantic Humility.''}
\emph{Philosophy and Phenomenological Research} 87 (2): 251--78. doi:
\href{https://doi.org/10.1111/j.1933-1592.2011.00550.x}{10.1111/j.1933-1592.2011.00550.x}.

\bibitem[\citeproctext]{ref-OppyRaivars2026}
Oppy, Graham, and Andrew Raivars. 2026. {``Blumberg and Holgu{ı́}n on
Unique Fictional Worlds.''} \emph{Philosophical Exchange}, 1--9. doi:
\href{https://doi.org/10.1080/29960126.2026.2712907}{10.1080/29960126.2026.2712907}.

\bibitem[\citeproctext]{ref-Schaffer2004}
Schaffer, Jonathan. 2004. {``Counterfactuals, Causal Independence and
Conceptual Circularity.''} \emph{Analysis} 64 (4): 299--309. doi:
\href{https://doi.org/10.1093/analys/64.4.299}{10.1093/analys/64.4.299}.

\bibitem[\citeproctext]{ref-Schulz2017}
Schulz, Moritz. 2017. \emph{Counterfactuals and Probability}. Oxford:
Oxford University Press.

\bibitem[\citeproctext]{ref-Stalnaker1968}
Stalnaker, Robert C. 1968. {``A Theory of Conditionals.''} In
\emph{Studies in Logical Theory}, edited by Nicholas Rescher, 98--112.
American Philosophical Quarterly Monograph Series 2. Oxford: Blackwell.

\bibitem[\citeproctext]{ref-Stalnaker1980}
---------. 1980. {``A Defense of Conditional Excluded Middle.''} In
\emph{Ifs: Conditionals, Belief, Decision, Chance, and Time}, edited by
William L. Harper, Robert Stalnaker, and Glenn Pearce, 87--104.
Dordrecht: Reidel. doi:
\href{https://doi.org/10.1007/978-94-009-9117-0_4}{10.1007/978-94-009-9117-0\_4}.

\bibitem[\citeproctext]{ref-Stein2019}
Stein, Megan. 2019. {``{NCIS: LA} Will Finally Reveal Who Won the Coin
Flip on the {`{JAG}'} Finale.''} \emph{Country Living}, May 12, 2019.
\url{https://www.countryliving.com/life/entertainment/a27396842/jag-who-won-coin-flip/}.

\bibitem[\citeproctext]{ref-Walton1990}
Walton, Kendall L. 1990. \emph{Mimesis as Make-Believe: On the
Foundations of the Representational Arts}. Cambridge, MA: Harvard
University Press.

\bibitem[\citeproctext]{ref-WilliamsWoodward2021}
Williams, J. Robert G., and Richard Woodward. 2021. {``The Cognitive
Role of Fictionality.''} \emph{Philosophy and Phenomenological Research}
102 (2): 423--38. doi:
\href{https://doi.org/10.1111/phpr.12660}{10.1111/phpr.12660}.

\bibitem[\citeproctext]{ref-Woodward2012}
Woodward, Richard. 2012. {``Fictionalism and Incompleteness.''}
\emph{No{û}s} 46 (4): 781--90. doi:
\href{https://doi.org/10.1111/j.1468-0068.2010.00826.x}{10.1111/j.1468-0068.2010.00826.x}.

\end{CSLReferences}




\end{document}
